\documentclass[a4paper,11pt]{article}

\usepackage{fontspec}
\setmainfont{Carlito}
\usepackage[margin=0.6in, top=0.5in, bottom=0.5in]{geometry}
\usepackage{titlesec}
\usepackage{enumitem}
\usepackage{xcolor}
\usepackage[colorlinks=true]{hyperref}
\usepackage{microtype}
\usepackage{ragged2e}
\usepackage{needspace}

% ---- palette: one accent, one ink, one muted grey ----
\definecolor{navy}{RGB}{12, 83, 166}
\definecolor{ink}{RGB}{33, 37, 41}
\definecolor{muted}{RGB}{96, 102, 110}
\hypersetup{urlcolor=navy, linkcolor=navy}
\color{ink}

% ---- ragged right, no hyphenation: even word spacing, no split words ----
\hyphenpenalty=10000 \exhyphenpenalty=10000 \emergencystretch=3em
\setlength{\parindent}{0pt}
\pagestyle{empty}
\linespread{1.03}
\AtBeginDocument{\RaggedRight}

% ---- type scale (ratio ~1.2): body 11 -> headline 12 -> section 13 -> name 22 ----
\titleformat{\section}
  {\color{navy}\fontsize{13}{15}\selectfont\bfseries\addfontfeature{LetterSpace=6}}
  {}{0em}{\MakeUppercase}[{\color{navy}\vspace{1pt}\titlerule[0.8pt]}]
\titlespacing*{\section}{0pt}{9pt}{5pt}

\setlist[itemize]{label={\color{navy}\small\textbullet}, leftmargin=1.1em, labelsep=0.45em,
  itemsep=1.5pt, parsep=0pt, topsep=2pt, partopsep=0pt}

% ---- building blocks: bold = organisation/project, italic = role/stack, grey = dates/places ----
\newcommand{\role}[4]{% org, location, dates, title
  \textbf{#1}{\color{muted}\,\textperiodcentered\,\textit{#2}}\hfill{\color{muted}#3}\\
  \textit{#4}\par}
\newcommand{\roleNoLoc}[3]{\textbf{#1}\hfill{\color{muted}#2}\\ \textit{#3}\par}
\newcommand{\project}[3]{% name, stack, right-hand tag
  \textbf{#1}\hfill#3\\ {\color{muted}\small\textit{#2}}\par}
\newcommand{\gh}[1]{\href{#1}{\small GitHub}}
\newcommand{\badge}[1]{{\color{navy}\small\textbf{#1}}}
\newcommand{\skillgroup}[2]{{\color{navy}\textbf{#1}}\\ #2\par\vspace{3pt}}
\newcommand{\gap}{\vspace{5pt}}
\newcommand{\num}[1]{\textbf{#1}}   % result numbers only


\begin{document}

% ================= HEADER =================
\begin{center}
  {\fontsize{22}{26}\selectfont\bfseries\color{navy} Thiruvel Andagurunathan Pandian}\\[4pt]
  {\fontsize{12}{14}\selectfont\color{navy} Data Scientist}\\[4pt]
  {\small
  Bristol, UK\nobreak\ \,\textperiodcentered\, \href{mailto:apthiruvel@gmail.com}{apthiruvel@gmail.com}\nobreak\ \,\textperiodcentered\, \href{tel:+447741016954}{+44 7741 016954}\nobreak\ \,\textperiodcentered\, \href{https://www.linkedin.com/in/thiruvel-a-p/}{linkedin.com/in/thiruvel-a-p}\nobreak\ \,\textperiodcentered\, \href{https://github.com/Thiruvel-AP}{github.com/Thiruvel-AP}\nobreak\ \,\textperiodcentered\, \href{https://thiruvel-ap.github.io/PersonalPortfolio/}{Portfolio}}\\[3pt]
  {\small\itshape\color{muted} Right to Work in UK\nobreak\ \,\textperiodcentered\, Available to join Immediately}

\end{center}

% ================= PROFILE =================
\needspace{5\baselineskip}\section{Profile}
Builds the modelling a brief needs and reports its limits plainly: gradient boosting (XGBoost churn), LSTM time series forecasting with time ordered validation, text clustering (SBERT, K-Means, DBSCAN) and 11,459 per gene models (train Pearson 0.807 vs test 0.545). Works with clients: scoped a loosely defined brief at 180 Degrees Consulting into a stakeholder demo (five star rating) and contributed to the client\textquotesingle{}s production MCP server, now in live use. Deploys on AWS and GCP (Lambda, Cloud Run) with CI/CD and MLflow experiment tracking. Owns data workflows end to end in Python and SQL, from merging structured, API and web scraped data through feature creation, tuning, validation and deployment (Docker, GitHub Actions CI/CD, MLflow, AWS Lambda).

% ================= SKILLS =================
\needspace{5\baselineskip}\section{Skills}
\skillgroup{Machine Learning, Deep Learning \& Statistics}{Predictive modelling (gradient boosting with XGBoost, random forest, scikit-learn), clustering (K-Means, HAC, DBSCAN, Leiden), deep learning (PyTorch, TensorFlow/Keras, LSTM, GRU, BiRNN, CNN), time series forecasting (cyclic time encoding, TimeSeriesSplit), anomaly detection (Isolation Forest), class imbalance handling (SMOTE), hyper parameter tuning (grid search, two stage Bayesian), model evaluation (held out, cross validation, Pearson, MCC), statistical inference and experimental design \textit{(confidence intervals, bootstrap, hypothesis testing; MSc coursework)}, MLflow experiment tracking}
\skillgroup{Python, Data Quality \& Cloud}{Python, pandas, NumPy, SciPy, Matplotlib, Seaborn, Plotly, SQL, PostgreSQL, R (tidymodels), identifier crosswalks, parquet, reproducible pipelines, AWS (Lambda, ECR, IAM, S3, Bedrock), GCP (Cloud Run), Docker, Kubernetes (AWS EKS), serverless deployment (Lambda, Cloud Run), containerisation, GitHub Actions CI/CD, FastAPI, TensorFlow Serving, pytest, Git}
\skillgroup{NLP, LLMs, Embeddings \& Evaluation}{NLP (text preprocessing, TF-IDF, SBERT embeddings, text clustering), Hugging Face Transformers, large language models (Claude, OpenAI, open weight via Ollama), prompt engineering, RAG with embeddings (ChromaDB), agents (LangChain, LangGraph), team LoRA fine tune of Llama, LLM evaluation (hallucination rate analysis, deterministic guardrails)}
\skillgroup{Client Delivery \& Communication}{Scoping ambiguous problems with clients, stakeholder demos, working with client engineering teams, code review, transparent reporting of model limitations, root cause analysis, technical documentation}

% ================= EXPERIENCE =================
\needspace{5\baselineskip}\section{Work Experience}
\needspace{4\baselineskip}
\role{AstraZeneca}{Bristol, UK (Remote)}{May 2026 – Sep 2026}{Data Scientist}
\begin{itemize}
  \item Diagnosed where it struggled with a three tier, \textasciitilde{}65 test evaluation harness (tool level → API → live) and hallucination rate analysis, debugging failures layer by layer, and used the same harness to validate the Groq to Bedrock migration after a rate limit evaluation.
  \item Built the Python data pipeline: reconciled \num{14 heterogeneous source datasets} through identifier crosswalks so the agent used trusted data, then deployed it as a Docker image on AWS Lambda (ECR, least privilege IAM) and documented the architecture and evaluation.
\end{itemize}
\gap
\needspace{4\baselineskip}
\role{180 Degrees Consulting Bristol}{Bristol, UK (Remote)}{Sep 2025 – Dec 2025}{Data Consultant (AI Engineering \& AI System)}
\begin{itemize}
  \item A product startup client needed an agentic AI prototype turned into software their team could use: led the architecture for a five person cross functional team, designed hierarchical agent workflows (MCP, Exograph GraphQL, PostgreSQL for persistent storage and retrieval across agent sessions) and built the Python prototype (FastMCP, Ollama, GitHub Actions CI), then ran a live demo with client stakeholders \num{(five star rating from senior leadership)} within the three month engagement.
  \item Carried it into production with the client\textquotesingle{}s engineering team: contributed to their production TypeScript/Node.js MCP server (OpenAI based query generation, Convex), now in live use, and reviewed pull requests in the client\textquotesingle{}s production repository alongside their CTO.
\end{itemize}
\gap
\needspace{4\baselineskip}
\role{Freelance}{Chennai, India}{Jan 2025 – Aug 2025}{AI/ML Engineer}
\begin{itemize}
  \item Built scikit-learn pipelines and LSTM/BiRNN/CNN models in TensorFlow/Keras, deployed as FastAPI services in Docker.
  \item Built a production RAG and agentic retrieval system (Claude API, LangChain, LangGraph, ChromaDB) served through a FastAPI REST API in Docker on GCP Cloud Run.
  \item Designed reproducible scikit-learn pipelines with feature engineering (temporal features, interaction terms, target encoding), model selection, cross validation and hyperparameter tuning across classification and regression tasks.
\end{itemize}
\gap
\needspace{4\baselineskip}
\role{Avasoft}{Chennai, India}{Jan 2023 – Sep 2024}{Software Development Engineer (Jun 2023 – Sep 2024); Trainee (Jan 2023 – May 2023)}
\begin{itemize}
  \item Delivered \num{10 reusable SwiftUI/UIKit component templates} that accelerated the team across concurrent releases, so features shipped from shared, tested building blocks.
\end{itemize}

% ================= PROJECTS =================
\needspace{5\baselineskip}\section{Projects}
\needspace{4\baselineskip}
\setbox0=\hbox{\textbf{Mining Insights from Customer Feedback: NLP \& unsupervised analytics (MSc team project)}\hspace{1em}\url{https://github.com/Thiruvel-AP/TextAnalytics}}%
\setbox2=\hbox{{\small\textit{Python, scikit-learn, SBERT, K-Means, HAC, DBSCAN}}\hspace{1em}{\small\url{https://github.com/Thiruvel-AP/TextAnalytics}}}%
\ifdim\wd0>\linewidth
\ifdim\wd2>\linewidth
\textbf{Mining Insights from Customer Feedback: NLP \& unsupervised analytics (MSc team project)}\\ {\color{muted}\small\textit{Python, scikit-learn, SBERT, K-Means, HAC, DBSCAN}}\\ \hfill{\small\url{https://github.com/Thiruvel-AP/TextAnalytics}}\par
\else
\textbf{Mining Insights from Customer Feedback: NLP \& unsupervised analytics (MSc team project)}\\ {\color{muted}\small\textit{Python, scikit-learn, SBERT, K-Means, HAC, DBSCAN}}\hfill{\small\url{https://github.com/Thiruvel-AP/TextAnalytics}}\par
\fi
\else
\project{Mining Insights from Customer Feedback: NLP \& unsupervised analytics (MSc team project)}{Python, scikit-learn, SBERT, K-Means, HAC, DBSCAN}{\url{https://github.com/Thiruvel-AP/TextAnalytics}}
\fi
\begin{itemize}
  \item Built a causal inference uplift analysis of a retail promotion in R during my MSc project, which showed the true effect on repeat purchases.
\end{itemize}
\gap
\needspace{4\baselineskip}
\setbox0=\hbox{\textbf{CoverFit: insurance benchmarking API (Bristol × UWE AI Hackathon)}\hspace{1em}\url{https://github.com/Thiruvel-AP/AIHackathon-Phoenix}}%
\setbox2=\hbox{{\small\textit{FastAPI, pandas, LiteLLM, Pydantic, pytest}}\hspace{1em}{\small\url{https://github.com/Thiruvel-AP/AIHackathon-Phoenix}}}%
\ifdim\wd0>\linewidth
\ifdim\wd2>\linewidth
\textbf{CoverFit: insurance benchmarking API (Bristol × UWE AI Hackathon)}\\ {\color{muted}\small\textit{FastAPI, pandas, LiteLLM, Pydantic, pytest}}\\ \hfill{\small\url{https://github.com/Thiruvel-AP/AIHackathon-Phoenix}}\par
\else
\textbf{CoverFit: insurance benchmarking API (Bristol × UWE AI Hackathon)}\\ {\color{muted}\small\textit{FastAPI, pandas, LiteLLM, Pydantic, pytest}}\hfill{\small\url{https://github.com/Thiruvel-AP/AIHackathon-Phoenix}}\par
\fi
\else
\project{CoverFit: insurance benchmarking API (Bristol × UWE AI Hackathon)}{FastAPI, pandas, LiteLLM, Pydantic, pytest}{\url{https://github.com/Thiruvel-AP/AIHackathon-Phoenix}}
\fi
\begin{itemize}
  \item An insurance sponsor needed a fast, low cost, explainable benchmark (hackathon brief): paired a deterministic matching engine (FastAPI, pandas, Pydantic) with a single LLM call (Gemini via LiteLLM) for the rationale, covered by pytest tests; responses arrive in \num{under 3 seconds at \textasciitilde{}£0.003 per query}.
\end{itemize}

% ================= EDUCATION =================
\needspace{5\baselineskip}\section{Education}
\needspace{3\baselineskip}
\roleNoLoc{University of Bristol}{Sep 2025 – Sep 2026}{MSc Data Science, Distinction \textit{(Expected)}}
{\color{muted}\small Large scale Data Engineering, Statistical Computing \& Empirical Methods, AI \& Text Analytics, Visual Analytics, Data Science Methods \& Practices, Technology, Innovation, Business \& Society}\par
\gap
\needspace{3\baselineskip}
\roleNoLoc{St. Joseph\textquotesingle{}s College of Engineering}{Aug 2019 – May 2023}{B.E. Electrical and Electronics Engineering, 91\%, First Class with Distinction}

% ================= CERTIFICATIONS & ACHIEVEMENTS =================
\needspace{5\baselineskip}\section{Certifications \& Achievements}
\begin{itemize}
  \item \textbf{Lloyd\textquotesingle{}s Banking Group Datathon, 3rd place}: Detected outliers in financial services data under \num{200:1 class imbalance}, comparing Isolation Forest against ensemble classifiers and tracking experiments in MLflow.
  \item Model Context Protocol: Advanced Topics, Anthropic, \textit{Jun 2026}\nobreak\ \,\textperiodcentered\, Introduction to Model Context Protocol, Anthropic, \textit{Jun 2026}\nobreak\ \,\textperiodcentered\, AWS Academy Cloud Foundations, \textit{Nov 2025}\nobreak\ \,\textperiodcentered\, 180 Degrees Consulting Data Science Consultant Certificate, \textit{Dec 2025}\nobreak\ \,\textperiodcentered\, Bristol PLUS Award, \textit{Mar 2026}
\end{itemize}

\end{document}